\documentclass[12pt]{article}
\usepackage{amsmath, amssymb}
\usepackage{geometry}
\usepackage{hyperref}
\usepackage{parskip}
\usepackage{microtype}
\usepackage{booktabs}
\usepackage{xcolor}
\geometry{margin=1.3in}

\title{\textbf{Channel Imbalance:\\
A Unified Mechanics Account of Mental Illness}}
\author{Joseph Shields}
\date{2026}

\begin{document}
\maketitle

\begin{abstract}
Current psychiatric classification describes mental illness in terms of
symptom clusters --- what the patient reports experiencing. Unified Mechanics
offers a structural account underneath the symptoms: each major condition
corresponds to a specific imbalance in the three channels of the nervous system.
Anxiety is boundary channel hyperactivation. Depression is exploration channel
suppression. PTSD is the boundary channel locked to a historical event. Addiction
is the matter channel hijacked by artificial familiarity. ADHD is exploration
channel dominance with impaired familiarity formation. OCD is the boundary unable
to resolve into familiarity. Autism is high boundary sensitivity combined with
deep and stable familiarity formation. Grief is the matter channel reaching for
an object whose boundary no longer responds. Each condition is described in terms
of its channel state, and existing treatments are reinterpreted through the same
framework --- not to replace clinical practice but to reveal why effective
treatments work and what they are actually doing to the system.
\end{abstract}

%% ─────────────────────────────────────────────────────────
\section{Foundations}
%% ─────────────────────────────────────────────────────────

The three channels of the nervous system are described fully in the companion
paper \textit{The Three Natures in Biology}. In brief:

\begin{itemize}
  \item \textbf{Boundary channel:} the site of crossing between states.
        Biologically: the synapse, the insular cortex, interoception. The felt
        sense of where the organism stands in relation to its environment.
  \item \textbf{Light channel:} signal propagating into new territory. Biologically:
        the action potential, the dopaminergic system, hippocampal encoding.
        Curiosity, novelty, exploration.
  \item \textbf{Matter channel:} the settling of a state into familiarity.
        Biologically: long-term potentiation, the basal ganglia, myelination.
        Habit, automaticity, comfort.
\end{itemize}

A healthy nervous system cycles continuously through all three. A boundary is
crossed (acknowledgement), a signal propagates and explores (exploration), a new
familiarity forms (familiarity). The cycle then begins again from a new position.

Mental illness, in this framework, is what happens when the cycle is disrupted.
One channel becomes dominant, or one channel cannot complete its function, or the
handoff between channels fails. The symptoms that psychiatry observes are the
downstream expression of these structural failures.

%% ─────────────────────────────────────────────────────────
\section{Anxiety: Boundary Channel Hyperactivation}
%% ─────────────────────────────────────────────────────────

In anxiety the boundary channel is running at excessive gain. The system
registers threat in stimuli that do not warrant it, treating ordinary crossings
as dangerous. The insular cortex reports persistent high-pressure boundary states.
The autonomic nervous system remains in a state of readiness that was designed
for brief crossings, not continuous occupation.

The matter channel cannot form because familiarity requires the boundary to
settle. If the boundary never settles, nothing becomes familiar. The world
remains perpetually uncrossed, and the organism remains perpetually at the edge.

The light channel is affected secondarily: exploration requires a degree of
boundary stability to feel safe. Without it, novelty feels like threat rather
than opportunity.

\textbf{Why treatments work:}
Anxiolytics (benzodiazepines, SSRIs over time) reduce boundary channel gain ---
lowering the insular cortex's sensitivity so that ordinary crossings are
processed as ordinary. Vagal nerve stimulation directly reduces autonomic
activation, bringing the boundary channel out of sustained alert. Meditation
works by training voluntary boundary-pressure reduction: the practitioner learns
to observe boundary states without amplifying them, allowing familiarity to form
underneath. Exposure therapy works by forcing controlled boundary crossings in
safe conditions until the system learns that the crossing resolves --- and
familiarity forms where fear had been.

%% ─────────────────────────────────────────────────────────
\section{Depression: Exploration Channel Suppression}
%% ─────────────────────────────────────────────────────────

In depression the exploration channel goes quiet. The dopaminergic system stops
responding to novelty. What was once new produces no signal. The light channel is
dim. Nothing reaches forward.

The matter channel continues to run --- the person moves through familiar
patterns, familiar thoughts, familiar days --- but without the exploration channel
to bring in new crossings, those patterns produce no reward. Familiarity without
novelty is not comfort. It is hollowness. The cycle has stalled at its final
stage with no way back to the beginning.

The boundary channel still operates but has nothing to cross into. A boundary
that opens onto silence, again and again, produces its own particular quality of
despair. The person can feel that something is wrong --- the boundary is
reporting --- but the light channel brings no answer.

\textbf{Why treatments work:}
Ketamine disrupts stuck matter-channel patterns directly and rapidly, forcing the
system out of entrenched familiarity before the exploration channel has been
restored. This is why ketamine works on treatment-resistant depression within
hours rather than weeks --- it breaks the deadlock rather than waiting for the
exploration channel to recover. SSRIs restore baseline dopaminergic sensitivity
over weeks, gradually returning the exploration channel to function. Exercise is
one of the most reliable antidepressants known because it forces exploration
channel activation through physical novelty and exertion --- the body exploring
its own capacity. Psychedelics produce massive, controlled exploration channel
activation, introducing crossings so numerous and unfamiliar that the stuck
matter patterns are disrupted and new familiarity must form. The therapeutic
window after a psychedelic session is the period in which the system is open to
forming new familiarity rather than defaulting to old patterns.

%% ─────────────────────────────────────────────────────────
\section{PTSD: The Boundary Locked to a Historical Crossing}
%% ─────────────────────────────────────────────────────────

In PTSD a boundary crossing of overwhelming intensity has failed to resolve. The
event was too large for the available familiarity to absorb. The boundary
opened, the system could not close it, and it has remained open ever since ---
not in the present but replayed from the past, as if the crossing is still
occurring.

The body does not distinguish between a real boundary crossing and a replayed
one. The insular cortex fires the same signals. The autonomic system responds
with the same urgency. The organism is locked at the moment of an event that has
already ended, returning to it whenever the boundary channel is triggered by
anything sufficiently similar to the original crossing.

The matter channel cannot form around the event because every approach to it
reopens the boundary before familiarity can settle. The person cannot move
through the crossing because the crossing will not complete.

\textbf{Why treatments work:}
EMDR (eye movement desensitisation and reprocessing) works by activating the
boundary crossing in a controlled setting while the bilateral stimulation reduces
the autonomic gain on the event. The crossing is re-entered at lower intensity,
again and again, until the system is able to move through it and familiarity
forms on the other side. Prolonged exposure therapy operates on the same
principle: systematic re-crossing of the boundary under safe conditions until the
event is no longer experienced as present. Propranolol, administered during
memory reconsolidation, reduces the autonomic intensity of the replay, making it
possible for the boundary to settle. All effective PTSD treatments are
functionally identical: they find a way to let the boundary crossing complete so
that familiarity can finally form.

%% ─────────────────────────────────────────────────────────
\section{Addiction: The Matter Channel Hijacked}
%% ─────────────────────────────────────────────────────────

The matter channel is designed to be reached through the full cycle: boundary
crossed, exploration undertaken, familiarity earned. This sequence is what makes
natural rewards --- food, connection, achievement, rest --- feel satisfying. The
familiarity is real because the crossing was real.

Addictive substances short-circuit this cycle. They deliver the matter channel's
reward --- the feeling of familiarity and resolution --- without requiring the
boundary to be crossed or the exploration to occur. Dopamine floods the system on
demand. Familiarity arrives without being earned.

The consequence is that the natural cycle atrophies. The boundary crossings of
ordinary life produce less and less reward because the matter channel has been
calibrated to the artificial intensity. Nothing that requires the full cycle feels
worth doing. The substance is the only thing that reaches the matter channel at
all. Familiarity has been hijacked: the organism is now deeply familiar with the
substance and increasingly unfamiliar with everything else.

\textbf{Why treatments work:}
Cold cessation is difficult because it removes the only functioning path to the
matter channel while the natural cycle has not yet recovered. Replacement
therapies (methadone, nicotine replacement) maintain minimal matter-channel
access while the natural cycle is slowly restored. The work of recovery is the
work of re-learning how to cross boundaries, explore, and earn familiarity
through the full sequence --- which is why connection, purpose, and structured
activity are consistently the most effective long-term supports. They restore
the cycle rather than substituting for it.

%% ─────────────────────────────────────────────────────────
\section{ADHD: Exploration Dominant, Familiarity Impaired}
%% ─────────────────────────────────────────────────────────

In ADHD the exploration channel is dominant and the threshold for novelty is
low. What is familiar becomes flat very quickly because the light channel moves
on before the matter channel has time to consolidate. The person is not
deficient in dopamine --- they are more sensitive to the absence of novelty.
The exploration channel demands constant new crossings. Anything already familiar
produces no signal, so attention moves elsewhere.

This is not a failure of will. It is a channel balance in which the light channel
outpaces the matter channel's ability to form. The cycle runs faster than it can
settle.

\textbf{Why treatments work:}
Stimulants are paradoxical: they calm ADHD rather than stimulating it further.
The reason is that they normalise the dopaminergic signal rather than amplifying
it. By raising baseline dopamine, they reduce the threshold gap between familiar
and novel, allowing the matter channel to hold a state long enough to become
genuinely familiar. The exploration channel no longer needs to move on so quickly
because the current state feels sufficiently rewarding. Structure and routine
serve the same purpose: they artificially extend the familiarity phase, giving
the matter channel time to consolidate what the exploration channel has found.

%% ─────────────────────────────────────────────────────────
\section{OCD: The Boundary That Cannot Resolve}
%% ─────────────────────────────────────────────────────────

In OCD the boundary crosses but familiarity does not form. The system expects
resolution after the crossing and does not receive it. So it crosses again. And
again. Not because the person wishes to but because the channel logic demands
that a boundary crossing produces a settling, and when the settling does not
come the only available response is to try again.

The compulsion is not the disorder. The compulsion is the system's attempt to
solve the disorder. It is the matter channel being reached for through an action
that used to resolve the boundary but no longer does. The boundary keeps
reopening before familiarity can stabilise.

\textbf{Why treatments work:}
Exposure and response prevention (ERP) is the most effective treatment for OCD.
It works by forcing the system to sit with the open boundary without performing
the compulsion. The discomfort is intense because the boundary channel is
reporting an unresolved state and demanding action. But if the compulsion is
withheld, something happens that the system had stopped believing was possible:
the boundary eventually settles on its own. Familiarity forms. The cycle
completes without the compulsion. Each successful exposure teaches the system
that resolution is available without the ritual --- and gradually the boundary
learns to settle.

%% ─────────────────────────────────────────────────────────
\section{Autism: High Boundary Sensitivity, Deep Familiarity}
%% ─────────────────────────────────────────────────────────

Autism is not a deficit. It is a specific channel balance: the boundary channel
is highly sensitive, and the matter channel forms deep, stable, and reliable
familiarity. These two features are not independent --- they are aspects of the
same underlying calibration.

High boundary sensitivity means that crossings which are unremarkable to most
nervous systems register with full force. Sensory experience is intense because
the boundary channel processes all input at high gain. Social situations are
demanding because they involve continuous rapid boundary crossings with
unpredictable timing. The world is always somewhat louder than expected.

But the matter channel, once it has formed familiarity, holds it with exceptional
stability. Interests become deep rather than broad. Routines provide genuine
relief because they are the matter channel operating in conditions of low boundary
pressure. The familiar is not boring --- it is genuinely satisfying in a way that
shallow familiarity is not.

What is labelled as rigidity is the matter channel protecting itself from
boundary overload. What is labelled as special interest is the light channel
going very deep in a direction where the boundary pressure is manageable. Neither
is pathological. They are the predictable outputs of a nervous system with this
particular channel balance.

%% ─────────────────────────────────────────────────────────
\section{Grief: Familiarity Without a Boundary}
%% ─────────────────────────────────────────────────────────

Grief is a channel state without a clinical category adequate to describe it,
because it is not an illness. It is the matter channel working exactly as it
should, in a situation that has become impossible.

When a person loses someone they love, the matter channel has formed deep
familiarity with that person --- with their presence, their voice, their
particular way of being in the world. That familiarity does not end with the
loss. The matter channel continues to reach for it. It expects the boundary to
respond, as it always has, and receives silence.

Grief is that silence. The matter channel running the familiarity pattern and
finding nothing to cross into. The boundary where someone used to be.

This is why grief is physical. Why it arrives in waves rather than steadily.
Why it can be triggered by unexpected things --- a song, a smell, a particular
quality of afternoon light. Each trigger activates the familiarity pattern and
the boundary reaches forward and finds nothing.

Grief does not end when the person stops thinking about the loss. It ends ---
partially, gradually, in patches --- when new familiarity forms that does not
require the absent boundary to respond. The lost person becomes part of the
matter channel in a different way: no longer as a present crossing, but as part
of the accumulated weight of what has made the system what it is.

%% ─────────────────────────────────────────────────────────
\section{Summary Table}
%% ─────────────────────────────────────────────────────────

\begin{center}
\begin{tabular}{lll}
\toprule
Condition & Channel State & Effective Treatment Logic \\
\midrule
Anxiety    & Boundary hyperactivated          & Reduce boundary gain; allow familiarity \\
Depression & Exploration suppressed           & Force exploration channel back online \\
PTSD       & Boundary locked to past event    & Complete the unfinished crossing \\
Addiction  & Matter channel hijacked          & Restore the full natural cycle \\
ADHD       & Exploration dominant             & Normalise novelty threshold \\
OCD        & Boundary cannot resolve          & Let boundary settle without compulsion \\
Autism     & High boundary sensitivity        & Accommodate; not a disorder to fix \\
Grief      & Matter reaching for absent boundary & Time; new familiarity forming \\
\bottomrule
\end{tabular}
\end{center}

The standard model of mental illness treats conditions as categories of symptom.
This framework treats them as states of a system. The difference matters
clinically: a system account predicts which treatments will work and why, suggests
what to try when first-line treatments fail, and identifies the underlying
mechanism rather than managing the surface expression.

A chemical imbalance is real. But it is the downstream expression of a channel
imbalance. The chemistry is how the channel state is implemented in tissue. The
channel state is what needs to be understood.

\end{document}
